% theory/revision/locality.tex -- Task 3 (G7-2): what in Section 4 is classical, what is new, and
% the shortened section.  Sources: locality-sources.md (retrieval log, quotations, instrument
% gaps); every quotation below is checked verbatim against sources_cache/ by check_quotes.py.
%
% WHAT WAS SEEN (primary text):  Dryden--Strohmaier, eq. (1) and refs [7],[8] (full);
%   Huber 1959, Huber II 1961 and its Nachtrag (full GDZ scans);  Wolpert, Bull. AMS 1977 (full);
%   Hejhal LNM 548: first two pages of each chapter only, plus the zbMATH review Zbl 0347.10018;
%   Buser 1992: Internet Archive full-text-search snippets only (no page numbers).
% INSTRUMENT GAPS (not seen; never read as absences):  Hejhal LNM 548 Ch. 3 interior (the
%   elliptic trace formula, its hypotheses on h, p. 351);  Iwaniec, Spectral Methods (Thm 10.2);
%   McKean, CPAM 25 (1972) and its correction;  Wolpert, Annals 1979.
%
% DECISIONS
% (a) Classical, cite:
%   * Signature locality (Thm 4.1).  For surfaces the heat expansion depends on the genus only:
%     this is McKean's use of the heat kernel in the trace formula (Buser, as quoted below, says
%     "Following McKean [1], we plug the heat kernel pM of M into the trace formula to prove
%     Huber's theorem."; McKean's own text NOT SEEN -- pinpoint to be added with library access).
%     For orbifolds it is immediate from the locality of the coefficients [Donnelly; DGGW Thm 4.8]
%     and homogeneity of H^2.  -> Demote to a Proposition with the structural proof (5 lines).
%   * The trace formula with elliptic terms (Thm 4.6) is classical: DS eq. (1) states it and cites
%     Hejhal LNM 548 and Iwaniec; zbMATH's review of Hejhal says of Ch. III: "Hier werden erstmals
%     auch Gruppen \(\Gamma\) mit kompaktem Quotienten zugelassen, welche elliptische Elemente
%     enthalten."  -> Present Thm 4.6 as a citation (DS eq. (1); Hejhal Ch. 3; Iwaniec).
%   * The mechanism "difference of heat traces = difference of hyperbolic terms, of size
%     e^{-l^2/4t}" is the standard Selberg/Huber/McKean mechanism; Huber compares length and
%     eigenvalue spectra (Huber 1959, Satz 7--8) by Dirichlet series, not heat kernels, and notes
%     (Huber II, p. 388, fn. 4) that his relation "auch aus der allgemeinen Spurformel von A.
%     Selberg ... hergeleitet werden" could be.  -> Say so; no novelty claim for the mechanism.
%   * Geodesic counting: for closed surfaces Buser gives an explicit bound depending on the genus
%     only, quoted below (snippet; lemma number from a forward reference "Lemma 6.6.4", not joined
%     to the statement in any snippet).  Huber gives asymptotics with implicit constants.
% (b) Lemma 4.7 (admissibility of h_t): NOT replaced by a citation.  The class of test functions
%   in Hejhal's elliptic trace formula could not be read, so it is not confirmed that the cited
%   theorem covers h_t "in the form needed".  Keep the minimal part: Lemma 4.7 stays, shortened to
%   the approximation argument (about 10 lines), introduced as standard.  ACTION FOR THE REWRITE
%   SESSION (library access): read Hejhal LNM 548, Ch. 3 (pp. 326--354, the theorem on p. 351 or
%   near it) and Iwaniec Thm 10.2.  If the hypothesis is "h even, holomorphic in
%   |Im r| <= 1/2 + eps, h(r) = O((1+|r|)^{-2-eps})", h_t qualifies and Lemma 4.7 becomes the
%   sentence in (S3) below marked [IF CONFIRMED]; the Weyl input of remark412.tex then also goes.
% (c) New: (i) the orbifold counting bound with explicit constant (Lemma 4.8; elementary, the
%   surface analogue is Buser's); (ii) the explicit small-t bound C(A,l,D) t^{-1/2} e^{-l^2/4t}
%   with its range (Thm 4.9(b)) and the a-posteriori form (Thm 4.9(c)); (iii) sharpness: the
%   leading asymptotic of the difference at the first length where the weighted length spectra
%   differ, and the fact that the factor t^{-1/2} cannot be dropped (Thm 4.10, Cor 4.11).  These
%   are an explicit packaging of the classical mechanism, and are to be presented as such.

% =====================================================================================
% SHORTENED SECTION 4 (target about 4 pp.).  Plan, with the exact wording of the new/classical
% sentences.  Theorem numbers refer to the current manuscript.
% =====================================================================================

\section{Heat does not hear the shape}\label{sec:locality}

% (S1) Opening paragraph -- replaces the current §4.1 heading text and Thm 4.1 statement.
Write $\Orb=\Gamma\backslash\mathbb H^2$ with $\Gamma\subset PSL(2,\R)$ a cocompact Fuchsian group. That
the heat expansion cannot see the moduli is classical: for closed hyperbolic surfaces it goes back to
McKean's use of the heat kernel in Selberg's trace formula \cite{mckean1972} (see also
\cite[Ch.~9]{buser1992}), and for orbifolds it follows from the locality of the heat invariants
\cite{donnelly1976}, \cite[Thm~4.8]{dggw2008}. We record it in the form used below.

\begin{proposition}[Signature locality]\label{thm:locality}
For every $j\ge1$, $c_j(\Orb)=\Phi_j(\sigma(\Orb))$ with $\Phi_1(g;m)=\frac12(2g-2+\sum_i(1-\frac1{m_i}))$ and
$\Phi_j(g;m)=\alpha_{j-1}\Phi_1(g;m)+\sum_ib_{j-2}(m_i)$ for $j\ge2$.
\end{proposition}

\begin{proof}
This is Proposition~\ref{prop:heatinput}. Structurally: the regular-stratum coefficients are
integrals of universal local invariants of the metric \cite[Def.~4.7]{dggw2008}, constant on a
hyperbolic orbifold, and each cone point contributes Donnelly's local terms for the rotations of its
isotropy group \cite[\S4.1]{dggw2008}, \cite{donnelly1976}, which depend only on the order since any
two cone points of the same order have isometric neighbourhoods.
\end{proof}

% (S2) §4.2 Moduli: keep Prop 4.2 (Thurston) and Prop 4.3 (rigidity, Troyanov) as they are; replace
% the proof of Prop 4.4 by the one-line argument of G7-20 (countable mapping-class orbits:
% Teich(O) = R^d with d > 0, the isometry classes of a signature are the orbits of the countable
% mapping-class group, and a countable union of countable sets is not R^d); keep Cor 4.5.

% (S3) §4.3 "Where the shape becomes audible".  Theorem 4.6 becomes a cited statement:

The shape is visible only beyond all orders, through the closed geodesics, by the classical
mechanism of Selberg, Huber and McKean \cite{huber1959,mckean1972,hejhal1976}: two orbifolds of the
same signature have the same identity and elliptic terms in the trace formula, and their heat traces
differ by the difference of the hyperbolic terms. We make this explicit.

% Theorem 4.6 statement unchanged; its proof becomes:
%   "This is the Selberg trace formula for cocompact Fuchsian groups with elliptic elements, in the
%    normalization of \cite[eq.~(1)]{drydenstrohmaier2009}, who take it from
%    \cite{hejhal1976,iwaniec2002}, applied to $h_t$ (Lemma~\ref{lem:admissible})."
% Lemma 4.7: keep, opening with "The following approximation argument is standard; for torsion-free
%   groups see \cite{marklof2011}."
%   [IF CONFIRMED in Hejhal/Iwaniec, replace Lemma 4.7 by: "The function $h_t$ is even, entire and
%    $O((1+|r|)^{-N})$ in every horizontal strip, so it belongs to the class of test functions of
%    \cite[Ch.~3, Thm ~X]{hejhal1976}, \cite[Thm~10.2]{iwaniec2002}."]
% Lemma 4.8: keep; add after it:
%   "For closed surfaces Buser gives a bound depending on the genus alone, at most
%    $(g-1)e^{L+6}$ closed geodesics of length at most $L$ that are not iterates of short ones
%    \cite[Lemma~6.6.4]{buser1992}; the diameter in Lemma~\ref{lem:counting} is the price of the
%    elementary argument, which also covers cone points."
% Theorems 4.9, 4.10, Corollary 4.11: keep (with thm12iii.tex), and introduce them by:

What is new in this section is quantitative: an explicit bound with explicit constants and range
(Theorem~\ref{thm:quantlocality}), and its sharpness, the leading term of the difference at the
first length where the two length spectra differ, which shows that the factor $t^{-1/2}$ cannot be
removed (Theorem~\ref{thm:sharp}, Corollary~\ref{cor:prefactor}).

% Remark 4.12: remark412.tex.  Remark 4.13: keep, shortened to its last two sentences, and cite
%   Huber for surfaces (equivalence of length and eigenvalue spectra, \cite[Satz~7--8]{huber1959})
%   next to \cite[Thm~1.1]{drydenstrohmaier2009} for orbisurfaces.

% (S4) Theorem 1.2 in the introduction: (i) becomes "Every heat invariant is a function of the
%   signature alone (classical; Proposition 4.1)"; (ii) unchanged; (iii) as in thm12iii.tex.
%   §1.1 "Prior work": add one sentence --
%   "That the heat expansion of a closed hyperbolic surface sees only its genus, and that the
%    moduli enter only through exponentially small geodesic terms, is classical
%    \cite{mckean1972,huber1959,buser1992}; the orbifold version is immediate from
%    \cite{dggw2008,drydenstrohmaier2009}. Our contribution in Section~\ref{sec:locality} is the
%    explicit constant and its sharpness."
%   and remove Thm 1.2 from the list of headline results in the abstract's novelty framing.

% NEW BIBLIOGRAPHY ENTRIES (verify every field against Crossref / the books before use):
%   huber1959  H. Huber, Zur analytischen Theorie hyperbolischer Raumformen und Bewegungsgruppen,
%              Math. Ann. 138 (1959) 1--26.  [seen]
%   hejhal1976 D. A. Hejhal, The Selberg Trace Formula for PSL(2,R), Vol. 1, LNM 548, Springer
%              1976.  [Ch. 3 interior NOT seen]
%   iwaniec2002 H. Iwaniec, Spectral Methods of Automorphic Forms, 2nd ed., GSM 53, AMS 2002. [NOT seen]
%   mckean1972 H. P. McKean, Selberg's trace formula as applied to a compact Riemann surface,
%              Comm. Pure Appl. Math. 25 (1972) 225--246.  [NOT seen]
%   buser1992  P. Buser, Geometry and Spectra of Compact Riemann Surfaces, Birkhauser 1992.
%              [snippets only; lemma number and page to be verified]
